Personalized route recommendation is usually evaluated with clean navigation logs, explicit user histories, or proprietary trajectory data. This paper advances a narrower open-data alternative: preference-aware route-candidate reranking from OpenStreetMap (OSM) features and public GPS-derived movement signals. We present a working system that generates OSM route candidates, extracts interpretable route features, constructs contextual pseudo-history profiles, and reranks candidates using profile, prompt, hybrid, or trajectory-derived vehicle profile scoring. The paper makes validity an explicit part of the method by evaluating trace reconstruction and route-candidate reranking separately, rather than treating public GPS traces as unquestioned user histories. On an exploratory Toronto OSM public-trace probe, timestamp coverage and GPS-to-route proximity are strong, while route-distance ratio identifies where reconstruction quality must still improve. The primary public benchmark path is a 100-query Porto taxi trajectory evaluation with OSM driving-route candidates, feature/path oracles, bootstrap confidence intervals, a learned feature-ranker baseline, paired bootstrap comparisons, and NASR-style path-overlap metrics. In the current Porto benchmark, a trajectory-derived vehicle profile improves ranking metrics over shortest-distance and a generic profile, especially Hit@3 and MRR, while a supervised prior-query feature ranker provides a stronger internal learned baseline on Hit@3 and mean feature distance. Path-level metrics still expose limitations: shortest-distance remains competitive on path F1/NDTW, and oracles show that map matching and candidate generation shape the achievable ceiling. The result is initial public benchmark evidence for a transparent, reproducible, interpretable OSM route-candidate reranking framework under realistic open-data constraints.
